% ---------------------------------------------------------------------------
% Verwendung und formale Muster
% ---------------------------------------------------------------------------
\addchap{\DEEN{Verwendungshinweis}{How to Use This Template}}

Diese Vorlage zur Erstellung von wissenschaftlichen Arbeiten an der
Hamburger Fern-Hochschule (HFH) basiert auf den Richtlinien zur Erstellung
wissenschaftlicher Arbeiten an der HFH (Fachbereiche Wirtschaft und Recht
sowie Technik) und wurde mit \LaTeX{} erstellt. Die Vorlage stellt ein Angebot
dar, das von Ihnen genutzt werden \textit{kann}. Als Editor wird Overleaf
empfohlen. Seine Verwendung ist aber nicht verpflichtend. Diese Vorlage
ist in Anlehnung an das \texttt{.docx}-Dokument aus dem Modul zum
wissenschaftlichen Arbeiten der HFH erstellt.

Am einfachsten öffnen Sie über den Link in der
\href{https://github.com/ilyaZar/hfh-latex-template}{README} eine eigene
Overleaf-Kopie. Alternativ können Sie das Repository herunterladen oder
klonen. Passen Sie anschließend den Metadatenblock am Anfang von
\texttt{main.tex} an und löschen oder überschreiben Sie alle Beispielinhalte,
die Sie nicht benötigen. Die PDF wird mit
\texttt{latexmk -pdf main.tex} beziehungsweise durch erneutes Kompilieren in
Overleaf vollständig aktualisiert.

Die Nutzung der Vorlage entbindet Sie nicht davon, die entsprechenden
\textbf{Studienbriefe zum wissenschaftlichen Arbeiten}, die sich mit der
Erstellung und Gestaltung wissenschaftlicher Arbeiten in Ihrem konkreten
Studiengang befassen, zu lesen und die dort gegebenen Hinweise (z.\,B. zu
qualitativen Anforderungen an wissenschaftliche Arbeiten, zum Zitieren, zum
Umgang mit Quellen und zu weiteren formalen Aspekten) in Ihren Arbeiten
umzusetzen.

Die Nutzung der Vorlage entbindet Sie auch nicht davon, zu überprüfen, ob
Sie die \textbf{aktuell geltenden} Richtlinien und Formatierungsregeln
der HFH in Ihrer Abschlussarbeit korrekt umsetzen.
